% GENERATED FILE. Edit canonical lessons, then run npm run generate:books.
% canonical-source-hash: fnv1a64:673b1e1f4ccbb2c5
% canonical-lessons: PA-C253-niyam, PA-C253-notis, PA-C253-katar, PA-C253-hatna, PA-C253-ishtihar

\chapter{Rules at the Gate}
\label{ch:pa-a2-rules-at-the-gate}
\hlchaptermodality{253}

\begin{chapteropening}
\textbf{By the end of this chapter:} \emph{I can understand a rule or a notice at an office gate.}

Complete the last lesson of chapter 253: I can understand a rule or a notice at an office gate.
\end{chapteropening}

\section[\texorpdfstring{\pa{ਨਿਯਮ} (niyam)}{ਨਿਯਮ (niyam)}]{\pa{ਨਿਯਮ} (niyam) — a rule}

\label{lesson:PA-C253-niyam}



\begin{quote}
Before the new one: say the Punjabi for to take out, then the Punjabi for perhaps.
\end{quote}

\subsection*{You'll want to know: \pa{ਨਿਯਮ}}
\textbf{\pa{ਨਿਯਮ}} — \emph{niyam} — \textquotedblleft{}a rule\textquotedblright{}.

\begin{quote}
\pa{ਬੈਂਕ} \pa{ਦਾ} \pa{ਨਿਯਮ} \pa{ਸਖ਼ਤ} \pa{ਹੈ।}
\end{quote}

\begin{quote}
\emph{baiṅk dā niyam sakht hai.}
\end{quote}

\begin{quote}
The bank's rule is strict.
\end{quote}

A masculine noun.

\begin{cousinweb}
\textbf{\pa{ਨਿਯਮ}} (\emph{niyam}) is taken over from Sanskrit \emph{niyama}, \textquotedblleft{}any fixed rule or law\textquotedblright{}, a tatsama (Monier-Williams).
\end{cousinweb}

\begin{grammarlens}[title={What you've built}]
One more word for reading notices and messages.
\end{grammarlens}

\subsection*{Your turn}
\begin{itemize}
\raggedright
  \item \emph{Say it:} \emph{niyam}
  \item \emph{Say it:} \emph{niyam}, once more
  \item \emph{Say it:} \emph{baiṅk dā niyam sakht hai.} — The bank's rule is strict.
  \item \emph{Recall:} say the Punjabi for to take out, then the Punjabi for perhaps, then say \emph{niyam} again
\end{itemize}

\begin{tcolorbox}[breakable,colback=teal!4,colframe=teal!35!black,title={Before you move on}]
What does \pa{ਨਿਯਮ} mean? (\textquotedblleft{}A rule\textquotedblright{}.) Say it once more.
\end{tcolorbox}
\section[\texorpdfstring{\pa{ਨੋਟਿਸ} (noṭis)}{ਨੋਟਿਸ (noṭis)}]{\pa{ਨੋਟਿਸ} (noṭis) — a notice}

\label{lesson:PA-C253-notis}



\begin{quote}
Before the new one: say the Punjabi for perhaps, then the Punjabi for a rule.
\end{quote}

\subsection*{You'll want to know: \pa{ਨੋਟਿਸ}}
\textbf{\pa{ਨੋਟਿਸ}} — \emph{noṭis} — \textquotedblleft{}a notice\textquotedblright{}.

\begin{quote}
\pa{ਦਫ਼ਤਰ} \pa{ਦੇ} \pa{ਦਰਵਾਜ਼ੇ} \pa{ਉੱਤੇ} \pa{ਨਵਾਂ} \pa{ਨੋਟਿਸ} \pa{ਹੈ।}
\end{quote}

\begin{quote}
\emph{daftar de darvāze utte navāṁ noṭis hai.}
\end{quote}

\begin{quote}
There is a new notice on the office door.
\end{quote}

A masculine noun.

\begin{cousinweb}
\textbf{\pa{ਨੋਟਿਸ}} (\emph{noṭis}) is borrowed from English \emph{notice}.
\end{cousinweb}

\begin{grammarlens}[title={What you've built}]
One more word for reading notices and messages.
\end{grammarlens}

\subsection*{Your turn}
\begin{itemize}
\raggedright
  \item \emph{Say it:} \emph{noṭis}
  \item \emph{Say it:} \emph{noṭis}, once more
  \item \emph{Say it:} \emph{daftar de darvāze utte navāṁ noṭis hai.} — There is a new notice on the office door.
  \item \emph{Recall:} say the Punjabi for perhaps, then the Punjabi for a rule, then say \emph{noṭis} again
  \item \emph{Write it:} copy \textbf{\pa{ਨੋਟਿਸ}} once, checking the whole word
\end{itemize}

\begin{tcolorbox}[breakable,colback=teal!4,colframe=teal!35!black,title={Before you move on}]
What does \pa{ਨੋਟਿਸ} mean? (\textquotedblleft{}A notice\textquotedblright{}.) Say it once more.
\end{tcolorbox}
\section[\texorpdfstring{\pa{ਕਤਾਰ} (katār)}{ਕਤਾਰ (katār)}]{\pa{ਕਤਾਰ} (katār) — a queue; a line of people}

\label{lesson:PA-C253-katar}



\begin{quote}
Before the new one: say the Punjabi for a rule, then the Punjabi for a notice.
\end{quote}

\subsection*{You'll want to know: \pa{ਕਤਾਰ}}
\textbf{\pa{ਕਤਾਰ}} — \emph{katār} — \textquotedblleft{}a queue; a line of people\textquotedblright{}.

\begin{quote}
\pa{ਬੈਂਕ} \pa{ਵਿੱਚ} \pa{ਲੰਮੀ} \pa{ਕਤਾਰ} \pa{ਹੈ।}
\end{quote}

\begin{quote}
\emph{baiṅk vich lammī katār hai.}
\end{quote}

\begin{quote}
There is a long queue in the bank.
\end{quote}

A feminine noun.

\begin{cousinweb}
\textbf{\pa{ਕਤਾਰ}} (\emph{katār}) is a loan, usually traced through Persian to Arabic \emph{qitār}, \textquotedblleft{}a line of camels\textquotedblright{}.
\end{cousinweb}

\begin{grammarlens}[title={What you've built}]
One more word for reading notices and messages.
\end{grammarlens}

\subsection*{Your turn}
\begin{itemize}
\raggedright
  \item \emph{Say it:} \emph{katār}
  \item \emph{Say it:} \emph{katār}, once more
  \item \emph{Say it:} \emph{baiṅk vich lammī katār hai.} — There is a long queue in the bank.
  \item \emph{Recall:} say the Punjabi for a rule, then the Punjabi for a notice, then say \emph{katār} again
  \item \emph{Write it:} copy \textbf{\pa{ਕਤਾਰ}} once, checking the whole word
\end{itemize}

\begin{tcolorbox}[breakable,colback=teal!4,colframe=teal!35!black,title={Before you move on}]
What does \pa{ਕਤਾਰ} mean? (\textquotedblleft{}A queue; a line of people\textquotedblright{}.) Say it once more.
\end{tcolorbox}
\section[\texorpdfstring{\pa{ਹਟਣਾ} (haṭṇā)}{ਹਟਣਾ (haṭṇā)}]{\pa{ਹਟਣਾ} (haṭṇā) — to move out of the way; to step back}

\label{lesson:PA-C253-hatna}



\begin{quote}
Before the new one: say the Punjabi for a notice, then the Punjabi for a queue.
\end{quote}

\subsection*{You'll want to know: \pa{ਹਟਣਾ}}
\textbf{\pa{ਹਟਣਾ}} — \emph{haṭṇā} — \textquotedblleft{}to move out of the way; to step back\textquotedblright{}.

\begin{quote}
\pa{ਲੋਕ} \pa{ਬੱਸ} \pa{ਲਈ} \pa{ਸੜਕ} \pa{ਤੋਂ} \pa{ਹਟਦੇ} \pa{ਹਨ।}
\end{quote}

\begin{quote}
\emph{lok bass laī saṛak ton haṭde han.}
\end{quote}

\begin{quote}
People move off the road for the bus.
\end{quote}

Past (I …): \pa{ਮੈਂ} \pa{ਸੜਕ} \pa{ਤੋਂ} \pa{ਹਟਿਆ}, \emph{maiṁ saṛak ton haṭiā}.

\begin{cousinweb}
\textbf{\pa{ਹਟਣਾ}} (\emph{haṭṇā}) is inherited from a reconstructed \emph{*haṭṭ}, \textquotedblleft{}move\textquotedblright{}; CDIAL lists Punjabi \emph{haṭṇā}, \textquotedblleft{}to move back or out of the way\textquotedblright{}, under it (CDIAL 13943).
\end{cousinweb}

\begin{grammarlens}[title={What you've built}]
One more word for reading notices and messages.
\end{grammarlens}

\subsection*{Your turn}
\begin{itemize}
\raggedright
  \item \emph{Say it:} \emph{haṭṇā}
  \item \emph{Say it:} \emph{haṭṇā}, once more
  \item \emph{Say it:} \emph{lok bass laī saṛak ton haṭde han.} — People move off the road for the bus.
  \item \emph{Recall:} say the Punjabi for a notice, then the Punjabi for a queue, then say \emph{haṭṇā} again
  \item \emph{Write it:} copy \textbf{\pa{ਹਟਣਾ}} once, checking the whole word
\end{itemize}

\begin{tcolorbox}[breakable,colback=teal!4,colframe=teal!35!black,title={Before you move on}]
What does \pa{ਹਟਣਾ} mean? (\textquotedblleft{}To move out of the way; to step back\textquotedblright{}.) Say it once more.
\end{tcolorbox}
\section[\texorpdfstring{\pa{ਇਸ਼ਤਿਹਾਰ} (ishtihār)}{ਇਸ਼ਤਿਹਾਰ (ishtihār)}]{\pa{ਇਸ਼ਤਿਹਾਰ} (ishtihār) — an advertisement}

\label{lesson:PA-C253-ishtihar}



\begin{quote}
Before the new one: say the Punjabi for a queue, then the Punjabi for to move out of the way.
\end{quote}

\subsection*{You'll want to know: \pa{ਇਸ਼ਤਿਹਾਰ}}
\textbf{\pa{ਇਸ਼ਤਿਹਾਰ}} — \emph{ishtihār} — \textquotedblleft{}an advertisement\textquotedblright{}.

\begin{quote}
\pa{ਅਖ਼ਬਾਰ} \pa{ਵਿੱਚ} \pa{ਨੌਕਰੀ} \pa{ਦਾ} \pa{ਇਸ਼ਤਿਹਾਰ} \pa{ਹੈ।}
\end{quote}

\begin{quote}
\emph{akhbār vich naukrī dā ishtihār hai.}
\end{quote}

\begin{quote}
There is a job advertisement in the newspaper.
\end{quote}

A masculine noun.

\begin{cousinweb}
\textbf{\pa{ਇਸ਼ਤਿਹਾਰ}} (\emph{ishtihār}) is a loan, usually traced through Persian to Arabic \emph{ishtihār}, \textquotedblleft{}making known\textquotedblright{}.
\end{cousinweb}

\begin{grammarlens}[title={What you've built}]
One more word for reading notices and messages.
\end{grammarlens}

\subsection*{Your turn}
\begin{itemize}
\raggedright
  \item \emph{Say it:} \emph{ishtihār}
  \item \emph{Say it:} \emph{ishtihār}, once more
  \item \emph{Say it:} \emph{akhbār vich naukrī dā ishtihār hai.} — There is a job advertisement in the newspaper.
  \item \emph{Recall:} say the Punjabi for a queue, then the Punjabi for to move out of the way, then say \emph{ishtihār} again
  \item \emph{Write it:} copy \textbf{\pa{ਇਸ਼ਤਿਹਾਰ}} once, checking the whole word
\end{itemize}

\begin{tcolorbox}[breakable,colback=teal!4,colframe=teal!35!black,title={Before you move on}]
What does \pa{ਇਸ਼ਤਿਹਾਰ} mean? (\textquotedblleft{}An advertisement\textquotedblright{}.) Say it once more.
\end{tcolorbox}
